\chapter{General Conclusion}
\label{chap:general_conclusion}

\section*{Conclusion}
\begin{flushleft}
This thesis explored the compelling intersection of \textbf{computer vision}, \textbf{deep reinforcement learning}, and \textbf{endoscopic surgery} to tackle one of the most technically demanding procedures in modern medicine: \textbf{colorectal polypectomy}. The human colon has an internal surface area of roughly 2,000\,cm\textsuperscript{2}. Navigating this extensive anatomical structure and accurately deploying a snare around a polyp extends far beyond a simple image classification task. It is a highly complex, sequential decision-making problem, as each localized movement fundamentally alters the visual observation space.

To solve this, we built a highly sophisticated, physics-driven \textbf{Unity simulation} to train vision-only agents. We split the procedure into two separate parts: \textbf{endoscope navigation} and \textbf{snare capture}. This decoupling significantly reduced the control complexity for the agent. The resulting policies demonstrated strong efficacy. In the first phase, our agent learned to safely steer the endoscope tip through those tortuous haustral folds, actively dodging collisions by reading real-time visual cues like \textbf{specular highlights} and dynamic shadows, achieving a robust \textbf{best score of 14.67} after 6.5 million steps. In the second phase, the agent mastered the highly precise task of looping a snare around a pedunculated polyp. This was facilitated by a carefully engineered \textbf{spatial reward system} and an \textbf{Intrinsic Curiosity Module (ICM)}. Both phases relied heavily on \textbf{Proximal Policy Optimization (PPO)} paired with \textbf{Convolutional Neural Networks (CNNs)} and \textbf{LSTM memory} to deliver remarkably safe and reliable policies.

Ultimately, the core contribution of this work is proving that we can build a software-first, structured foundation for \textbf{adaptive, semi-autonomous systems} that can assist doctors in highly dynamic operating environments. 
\end{flushleft}

\section{Limitations}

While we are confident in the results achieved in this project, there are naturally a few limitations to keep in mind. First, to ensure the simulation executes with high computational efficiency, we opted out of using true \textbf{soft-tissue deformation}. Instead, we modeled the environment using efficient \textbf{geometric sphere-proxies} rather than fully anatomical soft-body physics. Finally, taking this technology from the simulator to the clinic means bridging the \textbf{sim-to-real gap}. Given the timeframe of this thesis, we simply didn't have enough time or access to a physical endoscope to rigorously test our trained models in a real hospital setting.

\section{Future Work}

Looking ahead, there are several promising avenues to explore. For future work, it would be highly beneficial to integrate a high-end biomechanical engine like \textbf{SOFA (Simulation Open Framework Architecture)} to handle realistic tissue deformation. An additional enhancement would be the integration of a dedicated \textbf{object-detection AI} (like YOLO) to spot the polyps automatically. This would create a powerful \textbf{hybrid system} where the computer vision model provides target coordinates to the Deep RL agent. Bridging the sim-to-real gap with richer clinical datasets and eventually testing these intelligent policies on real-world \textbf{robotic endoscopy hardware} remains a long-term objective.
